\section{Pixel-to-parameter mapping}
\label{sec:mapping}

Before any fractal iteration can begin, every pixel on screen must be
assigned a complex parameter~$c$.  At moderate zoom depths this mapping is
a trivial affine transform, but at the extreme magnifications targeted by
\FractalShark{}---$10^{10{,}000}$ and beyond---the pixel-to-parameter
conversion becomes a precision-critical operation.  Adjacent pixels may
differ by $10^{-\text{thousands}}$, so rounding an absolute coordinate too early can
make neighboring pixels indistinguishable. The host constructs bounds and increments
at the selected arbitrary precision (\code{HighPrecision}). Kernel threads then
evaluate the affine map in their working type. Perturbation paths first subtract the
reference parameter at high precision, so the kernel receives small reference-relative
offsets instead of adding those offsets to a large rounded center.

This section describes how CUDA threads index pixels, how the coordinate
system is oriented, and how the affine map is parameterized.

Each CUDA thread computes a pixel coordinate \((X,Y)\) and maps it to a complex
parameter \(c = x_0 + i y_0\). A typical kernel starts with:
\begin{lstlisting}
int X = blockIdx.x * blockDim.x + threadIdx.x;
int Y = blockIdx.y * blockDim.y + threadIdx.y;
if (X >= width || Y >= height) return;
size_t idx = ConvertLocToIndex(X, height - Y - 1, width);
\end{lstlisting}

\subsection{Thread-to-pixel mapping}
A 2D CUDA grid of 2D blocks covers the image (\cref{fig:cuda-grid-mapping}).
Each thread is responsible for one pixel.
The bounds check prevents out-of-range threads from writing.

\begin{figure}[!htbp]
\centering
\begin{tikzpicture}[
    >=Stealth,
    blockcell/.style={draw, minimum width=0.65cm, minimum height=0.65cm,
                      inner sep=0pt},
    threadcell/.style={draw, minimum width=0.3cm, minimum height=0.3cm,
                       inner sep=0pt},
]

% ── CUDA grid: 3 x 2 blocks ──
\node[font=\scriptsize\bfseries] at (0.75, 1.85) {CUDA grid};

\foreach \bx in {0,1,2} {
  \foreach \by in {0,1} {
    \node[blockcell, fill=blue!8] (blk-\bx-\by)
      at (\bx*0.75, -\by*0.75) {};
  }
}
\node[blockcell, fill=orange!25, draw=orange!80!black, thick]
  at (blk-1-1) {};

\foreach \bx in {0,1,2} {
  \node[font=\tiny] at (\bx*0.75, 1.15) {\bx};
}
\node[font=\scriptsize] at (0.75, 1.45) {\code{blockIdx.x}};
\foreach \by in {0,1} {
  \node[font=\tiny] at (-0.55, -\by*0.75) {\by};
}
\node[font=\scriptsize, rotate=90, anchor=south] at (-0.85, -0.375)
  {\code{blockIdx.y}};

% ── Expanded block (1,1): 4 x 4 threads ──
\begin{scope}[shift={(0, -2.7)}]
  \node[font=\scriptsize\bfseries, anchor=south] at (0.54, 1.55)
    {Block $(1,\!1)$: $4\!\times\!4$ threads};

  \foreach \tx in {0,...,3} {
    \foreach \ty in {0,...,3} {
      \node[threadcell, fill=orange!8] (thr-\tx-\ty)
        at (\tx*0.36, -\ty*0.36) {};
    }
  }
  \node[threadcell, fill=orange!50, draw=orange!80!black, thick]
    at (thr-2-1) {};

  \foreach \tx in {0,...,3} {
    \node[font=\tiny] at (\tx*0.36, 1.05) {\tx};
  }
  \node[font=\tiny, anchor=south] at (0.54, 1.15)
    {\code{threadIdx.x}};
  \foreach \ty in {0,...,3} {
    \node[font=\tiny] at (-0.3, -\ty*0.36) {\ty};
  }
  \node[font=\tiny, rotate=90, anchor=south] at (-0.55, -0.54)
    {\code{threadIdx.y}};
\end{scope}

% Dashed arrow from grid block to expanded view
\draw[->, dashed, orange!60, thick]
  ([yshift=-2pt]blk-1-1.south) -- ++(0, -0.55);

% ── Formula ──
\node[draw, rounded corners=3pt, fill=gray!5, inner sep=5pt,
      font=\scriptsize, text width=4.0cm, align=left]
      (formula) at (5.0, -2.4) {
  $X = 1 \cdot 4 + 2 = 6$\\[2pt]
  $Y = 1 \cdot 4 + 1 = 5$\\[4pt]
  \centering Pixel $(6,\; 5)$
};

% Arrow from expanded thread to formula
\draw[->, thick] ([xshift=2pt]thr-2-1.east) -- (formula.west);

% ── Complex plane ──
\begin{scope}[shift={(9.2, -1.0)}]
  \node[font=\scriptsize\bfseries] at (1.0, 1.8) {Complex plane};

  \draw[->] (-0.3, -2.0) -- (2.4, -2.0)
    node[right, font=\scriptsize] {Re};
  \draw[->] (-0.2, -2.1) -- (-0.2, 1.5)
    node[above, font=\scriptsize] {Im};

  \draw[blue!25, fill=blue!3] (0.0, -1.7) rectangle (2.1, 1.2);
  \node[font=\tiny, blue!40] at (1.05, 1.35) {image};

  \fill[orange!80!black] (1.2, -0.1) circle (2.5pt);
  \node[font=\scriptsize, orange!80!black, right] at (1.35, -0.1)
    {$c = x_0 + iy_0$};

  \draw[dashed, gray!50, thin] (1.2, -0.1) -- (1.2, -2.0);
  \node[font=\tiny, gray, below] at (1.2, -2.05) {$x_0$};
  \draw[dashed, gray!50, thin] (1.2, -0.1) -- (-0.2, -0.1);
  \node[font=\tiny, gray, left] at (-0.25, -0.1) {$y_0$};
\end{scope}

% Arrow from formula to complex plane
\draw[->, thick] (formula.east) -- (8.85, -2.4)
  node[midway, above, font=\scriptsize] {affine map};

\end{tikzpicture}
\caption[CUDA thread-to-pixel-to-parameter mapping]{%
  CUDA thread-to-pixel-to-parameter mapping.
  A 2D grid of thread blocks (upper left) covers the image; each thread
  computes its pixel coordinate via
  \code{blockIdx}$\,\cdot\,$\code{blockDim}$\,+\,$\code{threadIdx}.
  The highlighted thread in block~$(1,1)$ at local index~$(2,1)$ yields
  pixel~$(6,5)$ (center), which an affine transform maps to complex
  parameter $c = x_0 + i\,y_0$ (right).}
\label{fig:cuda-grid-mapping}
\end{figure}

\subsection{Y-axis convention}
The expression \code{height - Y - 1} flips \(Y\). This flip is common when the image
buffer uses a top-left origin but the complex-plane mapping assumes a
bottom-left origin (or vice versa).

\subsection{Affine map into the complex plane}
The parameter \(c\) is computed by an affine transform:

\begin{align}
x_0 &= cx + dx \cdot X, \\
y_0 &= cy + dy \cdot Y,
\end{align}

where \code{cx,cy} anchor the plane (e.g., the coordinate at pixel \((0,0)\)),
and \code{dx,dy} are per-pixel increments. Different kernels compute these
expressions in different numeric types; the intent is always the same: map each
pixel to its associated complex parameter \(c\).

\paragraph{Worked viewport example.}
Take bounds $[-2,2]\times[-1,1]$ and an $800\times400$ output image with
supersampling factor $k=2$. The calculation grid has $1600\times800$ samples, with
$dx=4/1600=0.0025$ and $dy=2/800=0.0025$. In a bottom-left calculation coordinate
system, $(X,Y)=(800,400)$ gives $c=0+0i$. The corresponding top-left screen convention
uses $y=1-0.0025Y$ instead. Reversing a buffer row index changes storage order; it
must not introduce an additional reversal into the parameter calculation.

Near the real coordinate $1$, consecutive normal \code{float} values are separated
by $2^{-23}$. A step much smaller than that spacing can disappear when added to a
rounded absolute coordinate. The relevant spacing depends on the coordinate's
magnitude, whereas a reference-relative map avoids that large addition. HDR extends
the representable exponent range; the mantissa type still determines its precision.
